

How do normal cells become cancerous---and can we reverse their trajectory? How many genetic variants can we find in a large follow up datasets? 

Such questions exemplify the kind of scientific problems that artificial intelligence (AI) could, in principle, solve given ideal data and unlimited compute. In practice, however, scientific measurements are rarely ideal: data are scarce, noisy, and incomplete, and inference problems can be ill-posed given available observations. Unlike vision or language, where signals are abundant and noise relatively low, scientific data are limited in both size and signal. Yet many scientific domains come with deep prior knowledge about underlying mechanisms. My thesis will focuses on leveraging these sources of domain knowledge to make reliable inferences from imperfect data and ill posed problems. This perspective unifies my applications across diverse settings---from stochastic dynamics in biology to ocean current around a central question: \textit{\textbf{how can we design inference procedures when data are both limited in signal and structured by domain knowledge?}} My thesis will develop methods that embed domain-specific structure strongly enough to resolve ill-posedness, but flexibly enough to tolerate the imperfect or partial nature of these structures. My existing work on inferring trajectories using single cell data \cite{shen2025multi} was published in AISTATS 2025 with me as co-first author. Its follow up work \cite{berlinghieri2025oh} is under submission with me as co-first author. My existing work on predicting novel genetic variant counts \cite{shen2022double_trouble_journal} was accepted by \textit{Bayesian Analysis} with me as sole first author. For my planned future work, I have implemented the algorithms and am currently testing them.


%These are all scientific problems that might be solved by artificial intelligence with ideal data and endless compute. However, in practice, real-world measurements are rarely ideal. Data are often scarce or noisy, and inference problems can be fundamentally ill-posed given the observations available. But luckily we often have substantive domain knowledge to be used with these imperfect data or we can adaptively collect new, alternative form of data. Confronting these challenging inference problems leveraging these knowledge and data collection flexibility is central to my work across domains as varied as stochastic dynamics, human decision-making, astrophysics, and quantitative finance. At the core of my work lies a unifying question: \textit{\textbf{how can we design inference, learning, and data collection procedures when data are both limited in signal and structured by domain knowledge?}}

%These are all \textit{inverse} problems: we seek to recover hidden structure from indirect, noisy, or limited observations, a necessary step before designing interventions or controls.  

%My research tackles this question along two complementary directions address the inherent ill posedness and signal scarcity leveraging information from alternative sources. First, I develop methods for inference that efficiently leverage structural domain knowledge when they are available---solving the ill posedness by limiting on potential search spaces of models. Second, I study adaptive and sequential design---solving the ill posedness by strategically collect new, potentially alternative form of data of the system. Together, these directions advance probabilistic ML methodology while delivering practical tools that biologists and practitioners can use in their own domains.  

%I pursue this challenge along two complementary fronts. First \textbf{I develop methods that embed domain-specific structure strongly enough to resolve ill-posedness, but flexibly enough to tolerate the imperfect or partial nature of these structures}. By embedding structural domain knowledge directly into probabilistic machine learning, I turn ill-posed problems into tractable ones. Second, \textbf{I develop methods to collect the most informative new or alternative forms of data, and assess the robustness of conclusions drawn using current data}. These methods strategically collect information that fills in gaps left by scarce or noisy initial observations. Together, these strategies confront the fundamental limits of inference under scarcity, advancing probabilistic machine learning while equipping scientists with practical tools for their own domains.

%I pursue this challenge along two complementary fronts. First, \textbf{I develop methods that embed domain-specific structure strongly enough to resolve ill-posedness, yet flexibly enough to accommodate imperfect or partial knowledge.} By embedding structural priors directly into probabilistic machine learning, I turn ill-posed problems into tractable ones. Second, \textbf{I develop methods that both collect the most informative new data and assess the robustness of conclusions drawn from existing data.} These approaches strategically identify what information is most valuable---whether through new experiments or through critical evaluation of current evidence. Together, these strategies confront the fundamental limits of inference under data scarcity, advancing probabilistic machine learning while equipping scientists with reliable tools for their own domains.

%My thesis addresses not only the midstream of inferring stochastic dynamics under identifiability challenge but also upstream of the pipeline---from principled analysis to data collection. 

 %Upstream, my thesis will design adaptive data collection strategies that identify which new experiments or observations would most improve inference. 